\documentclass[12pt,aspectratio=1610,professionalfont,serif,t]{beamer}
\usepackage[osf,p]{newpxtext}
\usepackage{newpxmath}
\usepackage{fontawesome5}

\usepackage{fancyvrb}
\fvset{%frame=single,
commandchars=\\\{\},
%framesep=1mm,
fontfamily=helvetica,
fontsize=\normalsize
}
%\usepackage[bitstream-charter]{mathdesign}
\usepackage{listings}
\lstset{commentstyle=\color{orange}, keywordstyle=\color{yellow} }


\usetheme{Boadilla}
\useoutertheme{split}
\usecolortheme{albatross}


\setbeamertemplate{blocks}[rounded][shadow=false]
\setbeamersize{text margin left=\dimexpr 1cm+1ex\relax, text margin right=\dimexpr 1cm+1ex\relax}
\addtobeamertemplate{headline}{}{\usebeamerfont{frametitle}\vspace*{1ex}}
\setbeamertemplate{navigation symbols}{}
\setbeamertemplate{page number in head/foot}{}
%\beamerdefaultoverlayspecification{<+->}


\setbeamercolor*{structure}{fg=green!60!black,bg=black}
\setbeamercolor*{normal text}{fg=green!60!black,bg=black}
\setbeamercolor{palette primary}{use={structure,normal text},fg=green,bg=structure.bg}
\setbeamercolor{palette secondary}{use={structure,normal text},fg=structure.fg,bg=structure.bg}
\setbeamercolor{palette tertiary}{use={structure,normal text},fg=structure.fg,bg=structure.bg}
\setbeamercolor{palette quaternary}{use={structure,normal text},fg=white,bg=green!30!black}
\setbeamercolor{palette sidebar quaternary}{use={structure,normal text},fg=white,bg=green!30!black}
\setbeamercolor{title in head/foot}{use={structure,normal text},fg=white,bg=green!30!black}

\setbeamercolor*{example text}{use={structure},fg=green!60!black}
\setbeamercolor*{block body}{use={structure},bg=structure.bg}
\setbeamercolor*{block body alerted}{use={structure},bg=structure.bg}
\setbeamercolor*{block body example}{use={structure},bg=structure.bg}
\setbeamercolor*{block title}{parent=structure,bg=structure.bg}
\setbeamercolor*{block title alerted}{use={structure,alerted text},fg=alerted text.fg,bg=structure.bg}
\setbeamertemplate{navigation symbols}{}
\setbeamertemplate{items}[square]
\setbeamercolor{item projected}{fg=white}
\setbeamercolor*{normal text}{use={structure},fg=white!90!blue,bg=structure.bg}
\setbeamercolor*{separation line}{}
\setbeamercolor*{fine separation line}{}
\setbeamercolor{alerted text}{use={structure},fg=structure.fg}

\usepackage[italian]{babel}
\usepackage[utf8]{inputenc}
\usepackage{pgf}
\usepackage{verbatim}
\usepackage{inconsolata}
\usepackage{listings}
\lstset{language=C, frame=single,
	basicstyle=\ttfamily,
	numbers=left, numberstyle=\tiny\color{gray},
	numbersep=5pt, fancyvrb=true
}

\usepackage[noend]{algorithm2e}
\usepackage{pgf}
\usepackage{tikz}
\usepackage{graphicx}
\usetikzlibrary{%
	arrows,
	arrows.meta,
	positioning,
	calc,
	backgrounds,
	chains,
	matrix,
	patterns,
	automata,
	fit,
	graphs,
	decorations,
	decorations.pathmorphing,
	decorations.pathreplacing,
	decorations.markings,
}

\usepgflibrary{shapes,shapes.geometric}


\usepackage{url}


\usepackage{xmpmulti}
\usepackage{transparent}
\graphicspath{{img/}}

\newcommand{\zero}{\ensuremath{\mathtt{0}}\xspace}
\newcommand{\uno}{\ensuremath{\mathtt{1}}\xspace}
\newcommand{\due}{\ensuremath{\mathtt{2}}\xspace}
\newcommand{\tre}{\ensuremath{\mathtt{3}}\xspace}
\newcommand{\kleene}{\ensuremath{{}^{\ast}}\xspace}
\xspaceaddexceptions{\zero \uno \due \tre \kleene}


\newcommand{\simplebox}[1]{
      \begin{tikzpicture}
              \node[text centered, rectangle, draw=black, thick, text width=\linewidth,align=center, rounded corners, minimum height=1em] (C) {#1};
      \end{tikzpicture}
}
%
% In questo file si mettono i comandi utili solo per questo insegnamento.
%
\author{Gianluca Della Vedova}
\title{Linguaggi e Computabilità}
\institute{Univ. Milano--Bicocca\\
	\texttt{https://www.unimib.it/gianluca-della-vedova}}
\date{\today}



\begin{document}


\begin{frame}[fragile]
        \frametitle{Linguaggi e Computabilità: Lezione 20}
\begin{enumerate}
\item
      Forme normali
\item PDA deterministici e linguaggi prefix-free
\end{enumerate}
\end{frame}


\begin{frame}[fragile]
        \frametitle{Esercizio 1: forma normale di Chomsky}
Trasformare in forma normale di Chomsky la grammatica con variabile iniziale $S$ e produzioni
\begin{align*}
S & \to aAb \mid B \mid E\\
A & \to aA \mid \varepsilon\\
B & \to bB \mid b\\
E & \to cE
\end{align*}
Descrivere ogni passo: eliminazione delle $\varepsilon$-produzioni, delle produzioni unitarie,
dei simboli inutili, e infine riduzione dei corpi a due variabili o a un terminale.
\end{frame}
% SOLUZIONE
% 1. epsilon-produzioni. Annullabili: {A}. S non annullabile, quindi eps non in L.
%    S -> aAb | ab | B | E
%    A -> aA | a
%    B -> bB | b
%    E -> cE
% 2. Produzioni unitarie. Coppie unitarie: (S,S), (S,B), (S,E), (A,A), (B,B), (E,E).
%    S -> aAb | ab | bB | b | cE
%    A -> aA | a
%    B -> bB | b
%    E -> cE
% 3. Simboli inutili.
%    Generatori: a, b, c, A, B, S. E NON e' generatore (E -> cE non termina mai).
%    Elimino E e le produzioni che lo contengono:
%    S -> aAb | ab | bB | b ;  A -> aA | a ;  B -> bB | b
%    Raggiungibili da S: S, A, B (tutti) -> nessuna altra eliminazione.
% 4. Terminali nei corpi lunghi: Ca -> a, Cb -> b
%    S -> Ca A Cb | Ca Cb | Cb B | b
%    A -> Ca A | a
%    B -> Cb B | b
% 5. Corpi di lunghezza > 2: S -> Ca A Cb diventa S -> Ca D, D -> A Cb
% Grammatica finale in CNF:
%    S  -> Ca D | Ca Cb | Cb B | b
%    D  -> A Cb
%    A  -> Ca A | a
%    B  -> Cb B | b
%    Ca -> a
%    Cb -> b
% (L = { a^n b : n >= 1 } unione { b^m : m >= 1 })


\begin{frame}[fragile]
        \frametitle{Esercizio 2: forma normale di Chomsky}
Trasformare in forma normale di Chomsky la grammatica con variabile iniziale $S$ e produzioni
\begin{align*}
S & \to ASB \mid \varepsilon\\
A & \to aAS \mid a\\
B & \to SbS \mid A \mid bb
\end{align*}
Il linguaggio generato contiene $\varepsilon$: come si deve trattare?
\end{frame}
% SOLUZIONE
% 1. epsilon-produzioni. Annullabili: {S} (A e B non sono annullabili).
%    Poiche' eps in L(G), la CNF genera L(G) - {eps}
%    (se serve eps, si aggiunge una nuova variabile iniziale S0 -> S | eps).
%    S -> ASB | AB
%    A -> aAS | aA | a
%    B -> SbS | bS | Sb | b | A | bb
% 2. Produzioni unitarie: solo B -> A. Coppie unitarie: (B,A) oltre a quelle banali.
%    B -> SbS | bS | Sb | b | bb | aAS | aA | a
% 3. Simboli inutili. Generatori: A (a), B (b), S (AB). Tutti raggiungibili da S.
%    Nessuna eliminazione.
% 4. Terminali: Ca -> a, Cb -> b
%    S -> A S B | A B
%    A -> Ca A S | Ca A | a
%    B -> S Cb S | Cb S | S Cb | b | Cb Cb | Ca A S | Ca A | a
% 5. Corpi lunghi: introduco
%    D1 -> S B   (per A S B)
%    D2 -> A S   (per Ca A S, usato sia in A che in B)
%    D3 -> Cb S  (per S Cb S)
% Grammatica finale in CNF:
%    S  -> A D1 | A B
%    D1 -> S B
%    A  -> Ca D2 | Ca A | a
%    D2 -> A S
%    B  -> S D3 | Cb S | S Cb | b | Cb Cb | Ca D2 | Ca A | a
%    D3 -> Cb S
%    Ca -> a
%    Cb -> b


\begin{frame}[fragile]
        \frametitle{Esercizio 3: da PDA a DPDA}
Sia $P = (\{q_0, q_1\}, \{a,b\}, \{X, Z_0\}, \delta, q_0, Z_0)$ con accettazione per pila vuota e
\begin{align*}
\delta(q_0, a, Z_0) & = \{(q_0, XZ_0)\} &
\delta(q_0, a, X) & = \{(q_0, XX)\}\\
\delta(q_0, \varepsilon, X) & = \{(q_1, X)\} &
\delta(q_1, b, X) & = \{(q_1, \varepsilon)\}\\
\delta(q_1, \varepsilon, Z_0) & = \{(q_1, \varepsilon)\}
\end{align*}
\begin{enumerate}
\item dire perché $P$ non è deterministico e quale linguaggio $N(P)$ accetta;
\item verificare che $N(P)$ è prefix-free;
\item costruire un DPDA che accetta $N(P)$ per pila vuota;
\item si può fare lo stesso per $\{a^n b^m \mid n \ge m \ge 1\}$? Perché?
\end{enumerate}
\end{frame}
% SOLUZIONE
% 1. In q0 con X in cima alla pila sono possibili sia la mossa su a sia la
%    mossa eps verso q1: P deve "indovinare" quando sono finite le a.
%    N(P) = { a^n b^n : n >= 1 }.
% 2. Se a^m b^m fosse prefisso proprio di a^n b^n, con m < n, dopo a^m
%    dovrebbe esserci b, ma in a^n b^n c'e' a. Quindi nessuna stringa di L
%    e' prefisso proprio di un'altra stringa di L.
% 3. DPDA: si passa in q1 leggendo la prima b invece che con una mossa eps.
%    delta(q0, a, Z0)  = {(q0, X Z0)}
%    delta(q0, a, X)   = {(q0, XX)}
%    delta(q0, b, X)   = {(q1, eps)}
%    delta(q1, b, X)   = {(q1, eps)}
%    delta(q1, eps, Z0) = {(q1, eps)}
%    Deterministico: per ogni (stato, simbolo di pila) c'e' al piu' una mossa,
%    e quando c'e' una mossa eps (q1, Z0) non ci sono mosse su input.
%    Es. aabb: (q0,aabb,Z0) |- (q0,abb,XZ0) |- (q0,bb,XXZ0) |- (q1,b,XZ0)
%              |- (q1,eps,Z0) |- (q1,eps,eps)
% 4. No: {a^n b^m : n >= m >= 1} non e' prefix-free (aab e aabb stanno
%    entrambe in L). Un DPDA che accetta per pila vuota, dopo aver letto aab
%    ha la pila vuota e non puo' piu' muoversi, quindi non accetta aabb.
%    I linguaggi accettati da DPDA per pila vuota sono esattamente i DCFL
%    prefix-free. (Il linguaggio e' comunque accettato da un DPDA per stato
%    finale.)


\end{document}
